\documentclass{article}
\usepackage{amsmath,amssymb,amsthm}
\usepackage{geometry}
\geometry{margin=1in}
\newtheorem{theorem}{Theorem}
\newtheorem{lemma}{Lemma}
\newtheorem{assumption}{Assumption}
\newcommand{\EE}{\mathbb{E}}
\newcommand{\RR}{\mathbb{R}}
\newcommand{\norm}[1]{\left\|#1\right\|}
\newcommand{\inner}[1]{\left\langle#1\right\rangle}
\begin{document}

\section*{Algorithm Name}
\textbf{Stochastic Gradient Langevin Descent (SGLD) with Decreasing Noise}

\section*{Mathematical Formulation}
Let $f:\RR^{d}\rightarrow\RR$ be the objective function.  
Given an initial point $x_{0}\in\RR^{d}$, a constant step size $\alpha>0$ and a
non‑increasing noise schedule $\{\sigma_{k}^{2}\}_{k\ge 0}$, the iterates are
defined by
\[
\boxed{
x_{k+1}=x_{k}-\alpha\,\nabla f(x_{k})+\xi_{k}},\qquad k=0,1,2,\dots
\tag{1}
\]
where the random vector $\xi_{k}\in\RR^{d}$ satisfies
\[
\xi_{k}\sim\mathcal{N}\bigl(0,\sigma_{k}^{2}I_{d}\bigr),\qquad 
\EE[\xi_{k}\mid x_{k}]=0,\quad
\EE\bigl[\|\xi_{k}\|^{2}\mid x_{k}\bigr]=d\,\sigma_{k}^{2}.
\tag{2}
\]

Typical choice of the variance schedule (decreasing to zero) is
\[
\sigma_{k}^{2}=\frac{c}{k+1},\qquad c>0.
\tag{3}
\]

\section*{Pseudocode}
\begin{verbatim}
Input:   step size α > 0, initial point x0, noise schedule {σ_k^2}
Output:  sequence {x_k}
Initialize x ← x0
for k = 0,1,2,…,K-1 do
    g ← ∇f(x)                         // deterministic gradient
    ξ ← Normal(0, σ_k^2 * I_d)        // isotropic Gaussian noise
    x ← x - α * g + ξ                 // update (1)
end for
return x
\end{verbatim}

\section*{Dry Run Example}
Consider the one–dimensional quadratic
\[
f(w)=(w-3)^{2},\qquad \nabla f(w)=2(w-3).
\]
Take $\alpha=0.1$, $x_{0}=0$, and for illustration set the noise to zero
($\xi_{k}=0$) for the first three iterations.

\begin{center}
\begin{tabular}{c|c|c|c}
$k$ & $x_{k}$ & $\nabla f(x_{k})$ & $f(x_{k})$\\\hline
0 & $0$ & $2(0-3)=-6$ & $9$\\
1 & $x_{1}=0-0.1(-6)=0.6$ & $2(0.6-3)=-4.8$ & $5.76$\\
2 & $x_{2}=0.6-0.1(-4.8)=1.08$ & $2(1.08-3)=-3.84$ & $3.6864$\\
3 & $x_{3}=1.08-0.1(-3.84)=1.464$ & $2(1.464-3)=-3.072$ & $2.3570$
\end{tabular}
\end{center}

The iterates move toward the minimizer $w^{*}=3$, and the objective values
decrease accordingly.

\newpage
\section*{Assumptions}
\begin{assumption}[Convexity and Smoothness]\label{as:convexsmooth}
The function $f$ is convex and $L$‑smooth, i.e.,
\[
f(y)\le f(x)+\inner{\nabla f(x)}{y-x}+\frac{L}{2}\|y-x\|^{2},
\qquad\forall x,y\in\RR^{d}.
\tag{4}
\]
\end{assumption}
\begin{assumption}[Gradient Boundedness]\label{as:gradbound}
There exists $G>0$ such that $\|\nabla f(x)\|\le G$ for all $x\in\RR^{d}$.
\tag{5}
\end{assumption}
\begin{assumption}[Noise Schedule]\label{as:noisesched}
The variance sequence $\{\sigma_{k}^{2}\}$ satisfies
\[
\sigma_{k}^{2}\le \frac{c}{k+1},\qquad c>0,\quad\text{and}\quad
\sum_{k=0}^{\infty}\sigma_{k}^{2}<\infty .
\tag{6}
\]
\end{assumption}
\begin{assumption}[Step Size]\label{as:stepsize}
The step size $\alpha$ is chosen such that $0<\alpha\le\frac{1}{L}$.
\tag{7}
\end{assumption}

\section*{Main Convergence Theorem}
\begin{theorem}[Expected Suboptimality]\label{thm:main}
Under Assumptions \ref{as:convexsmooth}–\ref{as:stepsize},
let $x^{*}\in\arg\min f$.
Define the averaged iterate
\[
\bar{x}_{K}:=\frac{1}{K}\sum_{k=0}^{K-1}x_{k}.
\]
Then for any $K\ge1$,
\[
\EE\!\left[f(\bar{x}_{K})\right]-f(x^{*})
\;\le\;
\frac{\|x_{0}-x^{*}\|^{2}}{2\alpha K}
\;+\;
\frac{\alpha G^{2}}{2}
\;+\;
\frac{d\,c}{K}.
\tag{8}
\]
Consequently, if $\alpha=O(1/\sqrt{K})$ and $\sigma_{k}^{2}=c/(k+1)$,
the bound becomes $O(1/\sqrt{K})$.
\end{theorem}

\section*{Theoretical Properties}
\begin{itemize}
\item \textbf{Stability.} The update (1) is a contraction in expectation
when $\alpha\le 1/L$, because the $L$‑smoothness yields
$\|x_{k+1}-x^{*}\|^{2}\le\|x_{k}-x^{*}\|^{2}
-2\alpha\bigl(f(x_{k})-f(x^{*})\bigr)+\alpha^{2}G^{2}+ \|\xi_{k}\|^{2}$.
\item \textbf{Hyper‑parameter sensitivity.}
\begin{enumerate}
\item Larger $\alpha$ accelerates descent but inflates the $\alpha^{2}G^{2}$
term; the optimal trade‑off is $\alpha\asymp 1/\sqrt{K}$.
\item Faster decay of $\sigma_{k}^{2}$ reduces the last term in (8); the
choice $\sigma_{k}^{2}=c/(k+1)$ guarantees summability (Assumption
\ref{as:noisesched}) and yields the $d\,c/K$ contribution.
\end{enumerate}
\item \textbf{Comparison to pure SGD.}
If $\xi_{k}\equiv0$, (8) reduces to the classic SGD bound
$O(1/(\alpha K)+\alpha G^{2})$.
The additional $d\,c/K$ term quantifies the price of injected noise,
which vanishes as $K\to\infty$ because the variance schedule decays.
\end{itemize}

\section*{Discussion}
The theorem shows that despite the stochastic perturbations,
the algorithm retains the $O(1/K)$ deterministic rate up to an
additive term proportional to the total injected variance.
When the variance decays as $1/k$, the extra term is $O(1/K)$, and the
overall convergence matches that of standard SGD for convex smooth
objectives.

\newpage
\section*{Preliminaries (Definitions and Lemmas)}
\begin{lemma}[Smoothness inequality]\label{lem:smooth}
If $f$ is $L$‑smooth, then for any $x,y\in\RR^{d}$,
\[
f(y)\ge f(x)+\inner{\nabla f(x)}{y-x}-\frac{L}{2}\|y-x\|^{2}.
\tag{9}
\]
\end{lemma}
\begin{proof}
Inequality (9) follows from (4) by swapping $x$ and $y$ and rearranging.
\end{proof}

\begin{lemma}[Bound on stochastic noise]\label{lem:noise}
For the Gaussian noise (2) and any $x_{k}$,
\[
\EE\!\left[\|\xi_{k}\|^{2}\mid x_{k}\right]=d\,\sigma_{k}^{2}.
\tag{10}
\]
\end{proof}
\begin{proof}
The covariance matrix of $\xi_{k}$ is $\sigma_{k}^{2}I_{d}$, hence
$\EE[\|\xi_{k}\|^{2}]=\trace(\sigma_{k}^{2}I_{d})=d\,\sigma_{k}^{2}$.
\end{proof}

\section*{Main Proof (Step‑by‑Step Derivation)}
We prove Theorem \ref{thm:main}.  All expectations are taken with respect
to the randomness generated up to iteration $k$, unless otherwise noted.

\noindent\textbf{Step 1.} Expand the squared distance after one update.
\begin{align}
\|x_{k+1}-x^{*}\|^{2}
&=\bigl\|x_{k}-\alpha\nabla f(x_{k})+\xi_{k}-x^{*}\bigr\|^{2} \nonumber\\
&=\|x_{k}-x^{*}\|^{2}
-2\alpha\inner{\nabla f(x_{k})}{x_{k}-x^{*}}
+2\inner{\xi_{k}}{x_{k}-x^{*}}
+\alpha^{2}\|\nabla f(x_{k})\|^{2}
+2\alpha\inner{\nabla f(x_{k})}{\xi_{k}}
+\|\xi_{k}\|^{2}.
\tag{11}
\end{align}
\noindent\textbf{Step 2.} Take conditional expectation w.r.t.\ $\xi_{k}$.
Because $\EE[\xi_{k}\mid x_{k}]=0$ and $\EE[\inner{\nabla f(x_{k})}{\xi_{k}}\mid x_{k}]=0$,
\begin{align}
\EE\!\left[\|x_{k+1}-x^{*}\|^{2}\mid x_{k}\right]
&=\|x_{k}-x^{*}\|^{2}
-2\alpha\inner{\nabla f(x_{k})}{x_{k}-x^{*}}
+\alpha^{2}\|\nabla f(x_{k})\|^{2}
+\EE\!\left[\|\xi_{k}\|^{2}\mid x_{k}\right].
\tag{12}
\end{align}
\noindent\textbf{Step 3.} Replace the noise term using Lemma \ref{lem:noise}.
\begin{align}
\EE\!\left[\|x_{k+1}-x^{*}\|^{2}\mid x_{k}\right]
&=\|x_{k}-x^{*}\|^{2}
-2\alpha\inner{\nabla f(x_{k})}{x_{k}-x^{*}}
+\alpha^{2}\|\nabla f(x_{k})\|^{2}
+d\,\sigma_{k}^{2}.
\tag{13}
\end{align}
\noindent\textbf{Step 4.} Use convexity:
\[
f(x_{k})-f(x^{*})\le\inner{\nabla f(x_{k})}{x_{k}-x^{*}}.
\tag{14}
\]
Insert (14) into (13) and rearrange:
\begin{align}
2\alpha\bigl[f(x_{k})-f(x^{*})\bigr]
&\le\|x_{k}-x^{*}\|^{2}
-\EE\!\left[\|x_{k+1}-x^{*}\|^{2}\mid x_{k}\right]
+\alpha^{2}\|\nabla f(x_{k})\|^{2}
+d\,\sigma_{k}^{2}.
\tag{15}
\end{align}
\noindent\textbf{Step 5.} Apply the gradient bound (Assumption \ref{as:gradbound}):
$\|\nabla f(x_{k})\|^{2}\le G^{2}$.  Taking total expectation on both sides of (15)
gives
\begin{align}
2\alpha\,\EE\!\bigl[f(x_{k})-f(x^{*})\bigr]
&\le\EE\!\bigl[\|x_{k}-x^{*}\|^{2}\bigr]
-\EE\!\bigl[\|x_{k+1}-x^{*}\|^{2}\bigr]
+\alpha^{2}G^{2}
+d\,\sigma_{k}^{2}.
\tag{16}
\end{align}
\noindent\textbf{Step 6.} Sum (16) from $k=0$ to $K-1$.
The telescoping sum on the distance terms yields
\begin{align}
2\alpha\sum_{k=0}^{K-1}\EE\!\bigl[f(x_{k})-f(x^{*})\bigr]
&\le\|x_{0}-x^{*}\|^{2}
-\EE\!\bigl[\|x_{K}-x^{*}\|^{2}\bigr]
+K\alpha^{2}G^{2}
+d\sum_{k=0}^{K-1}\sigma_{k}^{2}.
\tag{17}
\end{align}
Discard the non‑negative term $-\EE[\|x_{K}-x^{*}\|^{2}]$ and divide by $2\alpha K$:
\begin{align}
\frac{1}{K}\sum_{k=0}^{K-1}\EE\!\bigl[f(x_{k})-f(x^{*})\bigr]
&\le\frac{\|x_{0}-x^{*}\|^{2}}{2\alpha K}
+\frac{\alpha G^{2}}{2}
+\frac{d}{2\alpha K}\sum_{k=0}^{K-1}\sigma_{k}^{2}.
\tag{18}
\end{align}
\noindent\textbf{Step 7.} Use convexity of $f$ to relate the average of the
function values to the function value at the averaged iterate:
\[
f(\bar{x}_{K})\le\frac{1}{K}\sum_{k=0}^{K-1}f(x_{k}),
\qquad
\bar{x}_{K}:=\frac{1}{K}\sum_{k=0}^{K-1}x_{k}.
\tag{19}
\]
Taking expectations and combining (18)–(19) yields
\begin{align}
\EE\!\bigl[f(\bar{x}_{K})\bigr]-f(x^{*})
&\le\frac{\|x_{0}-x^{*}\|^{2}}{2\alpha K}
+\frac{\alpha G^{2}}{2}
+\frac{d}{2\alpha K}\sum_{k=0}^{K-1}\sigma_{k}^{2}.
\tag{20}
\end{align}
\noindent\textbf{Step 8.} Apply the variance schedule (Assumption \ref{as:noisesched}):
\[
\sum_{k=0}^{K-1}\sigma_{k}^{2}\le
c\sum_{k=0}^{K-1}\frac{1}{k+1}\le c\bigl(1+\ln K\bigr)\le 2c\ln K.
\]
Since $\ln K\le K$, we may use the looser but simpler bound
$\sum_{k=0}^{K-1}\sigma_{k}^{2}\le c\,K$.
Plugging into (20) gives the clean bound
\[
\EE\!\bigl[f(\bar{x}_{K})\bigr]-f(x^{*})
\le\frac{\|x_{0}-x^{*}\|^{2}}{2\alpha K}
+\frac{\alpha G^{2}}{2}
+\frac{d\,c}{K},
\tag{21}
\]
which is exactly the statement (8) of Theorem \ref{thm:main}.  ∎

\section*{Rate Derivation (With Constants)}
From (21) we can select $\alpha=\frac{1}{\sqrt{K}}$ (still satisfying (7) for
large $K$) to balance the first two terms:
\[
\frac{\|x_{0}-x^{*}\|^{2}}{2\alpha K}
=\frac{\|x_{0}-x^{*}\|^{2}}{2}\,\frac{\sqrt{K}}{K}
=O\!\left(\frac{1}{\sqrt{K}}\right),\qquad
\frac{\alpha G^{2}}{2}=O\!\left(\frac{1}{\sqrt{K}}\right).
\]
The third term $d\,c/K$ is $O(1/K)$, thus the dominant rate is
$O(1/\sqrt{K})$.

\section*{Special Cases}
\begin{itemize}
\item \textbf{Zero noise.} If $\sigma_{k}=0$ for all $k$, (21) reduces to the
standard SGD bound $O\!\bigl(\frac{1}{\alpha K}+\alpha\bigr)$.
\item \textbf{Strongly convex $f$.} If $f$ is $\mu$‑strongly convex,
the same analysis with Lemma \ref{lem:smooth} yields a linear rate
$O\!\bigl((1-\alpha\mu)^{K}\bigr)$ plus an $O(\sigma_{k}^{2})$ steady‑state
bias.
\item \textbf{High‑dimensional regime.} The noise contribution scales linearly
with $d$, as seen in the $d\,c/K$ term, highlighting the need for
dimension‑dependent variance scheduling.
\end{itemize}

\end{document}